\documentclass[10pt,a4paper]{article}
\usepackage[margin=2.1cm]{geometry}
\usepackage{amsmath,amssymb}
\usepackage{booktabs,tabularx,array}
\usepackage[T1]{fontenc}
\usepackage{lmodern}
\usepackage{microtype}
\usepackage{xcolor}
\usepackage[colorlinks=true,linkcolor=blue!50!black,urlcolor=blue!50!black,citecolor=blue!50!black]{hyperref}
\usepackage{enumitem}
\setlist{nosep,leftmargin=1.4em}
\newcommand{\code}[1]{\texttt{\small #1}}
\newcommand{\dd}{\,\mathrm{d}}
\newcommand{\sinc}{\operatorname{sinc}}
\newcommand{\std}{\operatorname{std}}
\newcommand{\med}{\operatorname{median}}
\newcommand{\EI}{\operatorname{EI}}
\newcommand{\one}{\mathbf 1}
\newcommand{\caveat}[1]{\par\smallskip\noindent\colorbox{orange!10}{\parbox{\dimexpr\linewidth-2\fboxsep}{\small\textbf{Assumption / status.} #1}}\par\smallskip}
\newcommand{\keybox}[1]{\par\smallskip\noindent\colorbox{blue!6}{\parbox{\dimexpr\linewidth-2\fboxsep}{\small #1}}\par\smallskip}

\title{The models, in full:\\ trace physics features, the noise forest, the usability classifier,\\ the fleet Gaussian process and the Bayesian optimisation}
\author{ccd\_opt\_ml level-1 study (PA03 board 101, v3 night)}
\date{20 November 2026; forest on every input type, fleet-GP training and the search loop written out, the same day}
\sloppy

\begin{document}
\maketitle

\noindent This note writes down the mathematics of the models that are used to choose the CCD voltages, exactly as the
code evaluates them, and how each one is trained. The chain is:
\begin{center}\small
trace $\rightarrow$ physics features (Sec.~\ref{sec:features}) $\rightarrow$ noise forest (Sec.~\ref{sec:noise})
\quad and \quad
settings + trace $\rightarrow$ usability classifier (Sec.~\ref{sec:gate})\\[2pt]
images of other amplifiers $\rightarrow$ fleet Gaussian process (Secs.~\ref{sec:gp}--\ref{sec:fleet})
$\rightarrow$ expected improvement with a hard usability gate (Sec.~\ref{sec:bo})
\end{center}
Appendix~\ref{app:other} describes every other model that was tried on the same data, and which was best for each task.
Script names are in \code{monospace}, with line numbers at the commit of this note (branch
\code{v3-analysis-amp-physics}). Standard methods are cited; where no reference is given the step is our own. Shaded boxes
mark assumptions, approximations and places where an earlier document and the code disagree. The trace-to-physics step
is derived in more detail in the companion note \emph{From the readout trace to the amplifier physics}
(\code{docs/overleaf\_trace\_math\_2026-11-17}).

\tableofcontents

%=====================================================================================================================
\section{Data, labels and notation}\label{sec:data}

\paragraph{Data.} The v3 night on PA03 board 101, amplifiers (channels) 0, 1 and 3: 381 images (packets $\times$
amplifiers) over two sessions. \textbf{Session G}: 80 Sobol-drawn voltage settings, 14 nominal packets, 4 replicates.
\textbf{Session H}: 24 new settings and 5 nominal packets, scored after the scorecard's model choice was frozen
(\code{v3\_report.tex}:56--63; \code{lvl1/freeze\_choice\_v3.py}). Each packet takes, at one setting, three or four
0.2\,s raw traces per amplifier and one skipper image of 500 non-destructive reads per pixel \cite{tiffenberg2017}.

\paragraph{Labels} (from the WADERs peak fit of each image, \code{src/process\_packet.py}):
$\sigma_e$ the read noise in electrons after 500 skips, $\lambda$ the charge per pixel (dark current + clock-induced charge)
in e$^-$/pixel, $g$ the gain in ADU/e$^-$. Every model works in natural logarithms,
\begin{equation}
y_\sigma=\ln\sigma_e,\qquad y_\lambda=\ln\lambda,\qquad y_g=\ln g .
\end{equation}
An error $\hat y-y$ in log units, multiplied by 100, is quoted as ``\%'' throughout (for small errors it is the relative error).

\paragraph{Settings.} $\theta\in\mathbb R^{p}$ the voltages of one packet (the \code{theta.*} columns): the
high and low rails of SW, RG, V1--3, H1--3, plus VSUB, VDD and VR (two pins per rail where the board has two).
Different models use different encodings of $\theta$; each section says which.

\begin{center}\small
\begin{tabularx}{\linewidth}{@{}lX@{}}\toprule
symbol & meaning\\\midrule
$a$ & amplifier (here 101:0, 101:1, 101:3)\\
$x$ & a setting, as a standardised vector (the knobs $z$ of Sec.~\ref{sec:knobs})\\
$\phi$ & the trace physics features of one packet and amplifier (Sec.~\ref{sec:features})\\
$J$ & the objective, $J=w\,y_\sigma+(1-w)\,y_\lambda$ with $w\in\{1,\tfrac12,0\}$ (Sec.~\ref{sec:objective})\\
$f_s$ & ADC sample rate, $15\times10^6$\,s$^{-1}$\\
$\mathcal N(\mu,\sigma^2)$, $\Phi$, $\varphi$ & normal distribution, its cumulative distribution and density\\\bottomrule
\end{tabularx}
\end{center}

%=====================================================================================================================
\section{Trace physics features}\label{sec:features}

A trace is one channel's raw ADC samples $v[n]$ (ADU) with the sequencer flag bits $b_k[n]$, about 188\,ms at $f_s$.
The features are computed by \code{amp\_phys/trap\_from\_trace.py} and \code{amp\_phys/physics\_features.py}
(single-skip noise and the window fit) and \code{amp\_phys/gain\_features.py} (deterministic signal sizes).

\subsection{Skips and software CDS}
\begin{itemize}
\item Pedestal-window starts $e_k$: rising edges of flag bit 0. Skip period $P=\med_k(e_{k+1}-e_k)$ samples,
  $T=P/f_s$ (\code{trap\_from\_trace.py}:46--47).
\item Window length $w_0$: median length of the bit-0-high runs. Window spacing $s_0$: median offset from a pedestal start
  to the next rising edge of bit 1 (signal window); $\Delta t=s_0/f_s$ (lines 51--53).
\item Consecutive skips $P\pm3$ samples apart belong to one pixel; the first 20 skips of each pixel are dropped
  (lines 56--62; \code{physics\_features.py}:43).
\end{itemize}
For each kept skip $k$ and each software window length $\tau\in\{30,45,60,90,120,140\}$ samples
($2.0, 3.0, 4.0, 6.0, 8.0, 9.3$\,\textmu s), the CDS value averages the \emph{last} $\tau$ samples of each hardware window
(lines 65--69):
\begin{equation}
x_k(\tau)=\frac1\tau\sum_{j\in W^{\rm sig}_k(\tau)}v_j\;-\;\frac1\tau\sum_{j\in W^{\rm ped}_k(\tau)}v_j,
\end{equation}
with $W^{\rm ped}_k(\tau)=\{e_k+w_0-\tau,\dots,e_k+w_0-1\}$ and $W^{\rm sig}_k(\tau)=W^{\rm ped}_k(\tau)+s_0$.

\subsection{Single-skip noise $\sigma_1(\tau)$}
The charge in a pixel is the same for every skip, so the difference of consecutive skips of one pixel contains no signal:
$x_k=gQ+n_k\Rightarrow x_{k+1}-x_k=n_{k+1}-n_k$. Hence
\begin{equation}
\sigma_1(\tau)=\frac{1}{\sqrt2}\,\std_k\big(x_{k+1}(\tau)-x_k(\tau)\big)\qquad\text{[ADU]},
\end{equation}
over same-pixel pairs only (\code{trap\_from\_trace.py}:69).

\subsection{The window fit: white level $S_w$ and 1/f level $K$}\label{sec:windowfit}
For a stationary noise with one-sided power spectral density $S(f)$ (ADU$^2$/Hz), the variance of the skip-difference
estimator is $\sigma_1^2(\tau)=\int_0^\infty S(f)\,|H_\tau(f)|^2\dd f$ with the transfer function of two box windows
$\Delta t$ apart (CDS \cite{hopkinson1982}) times the skip difference divided by $\sqrt2$:
\begin{equation}
|H_\tau(f)|^2=\underbrace{4\sin^2(\pi f\Delta t)\,\sinc^2(f\tau)}_{\text{CDS}}\;\cdot\;\underbrace{2\sin^2(\pi fT)}_{\text{skip difference}/\sqrt2},
\qquad \sinc(u)=\frac{\sin\pi u}{\pi u}
\end{equation}
(\code{trap\_from\_trace.py}:37--42; no front-end roll-off, \code{FILT["fc"]}\,=\,None). The noise model is white plus 1/f,
\begin{equation}
S(f)=S_w+\frac{K}{f}\quad\Longrightarrow\quad
\sigma_1^2(\tau)=S_w\,I_w(\tau)+K\,I_f(\tau),\qquad I_w=\int|H_\tau|^2\dd f,\quad I_f=\int\frac{|H_\tau|^2}{f}\dd f,
\label{eq:windowfit}
\end{equation}
with the integrals as Riemann sums on 200\,000 log-spaced frequencies from 1\,Hz to 7.5\,MHz. $I_w\propto1/\tau$ while $I_f$
barely depends on $\tau$, so the six window lengths separate the two levels. $(S_w,K)$ is the weighted linear
least-squares solution that minimises the \emph{relative} residuals,
\begin{equation}
(S_w,K)=\arg\min\sum_{i=1}^{6}\left[\frac{S_wI_w(\tau_i)+KI_f(\tau_i)-\sigma_1^2(\tau_i)}{\sigma_1^2(\tau_i)}\right]^2
\end{equation}
(\code{np.linalg.lstsq} with weights $1/y_i$). Units: $S_w$ in ADU$^2$/Hz, $K$ in ADU$^2$.

\caveat{The fit is unconstrained: $S_w$ can come out negative (min $-0.021$ in the v3 table); the noise model drops those
rows. $K$ absorbs everything that does not average down with $\tau$: transistor 1/f, setup pickup and slow wander.}

\subsection{Predicted noise of a 500-skip pixel, $\sigma_{500}$}
The image averages $N=500$ skips. Its filter replaces the skip difference by the $N$-skip average:
\begin{equation}
|H^{(N)}(f)|^2=4\sin^2(\pi f\Delta t)\,\sinc^2(f\tau_{\rm img})\left[\frac{\sin(N\pi fT)}{N\sin(\pi fT)}\right]^2,
\qquad N=500,\ \tau_{\rm img}=10\,\text{\textmu s},
\end{equation}
on 300\,000 log-spaced frequencies from 0.1\,Hz (\code{physics\_features.py}:32--38). With the fitted levels,
\begin{equation}
\sigma_{500}=\sqrt{\max\!\big(S_w I_w^{(N)}+\max(K,0)\,I_f^{(N)},\,0\big)}\quad\text{[ADU]},
\label{eq:s500}
\end{equation}
and a white-only version $\sigma_{500}^{\rm white}=\sqrt{S_wI_w^{(N)}}$ (lines 49--52). $\sigma_{500}$ is the
noise a 500-skip pixel \emph{would} have in ADU; it is not divided by the gain here.

\caveat{The v3 feature table (\code{amp\_phys/results/per\_channel/physics\_features\_v3.csv}) was made at commit
\code{a532e6b} with $\tau_{\rm img}$ fixed at 10\,\textmu s. The current code uses the packet's own window; on v3 the timing
is nominal, so the two agree.}

\subsection{Deterministic signal sizes (\code{gain\_features.py})}\label{sec:gainfeat}
These are not noise; they are the sizes of clocked signals, each a known charge times the conversion gain. They are used
by the usability classifier and the trace-driven search (Sec.~\ref{sec:gate}), not by the 8.6\,\% noise forest.
\begin{itemize}
\item \textbf{Reset jump}: in the median horizontal-transfer waveform $\bar w$ (300 samples before to 1700 after each bit-6 edge,
  baseline-subtracted), the edge found about 278 samples before the H flag; jump $=\overline{\bar w}_{e+3..e+11}-\overline{\bar w}_{e-12..e-4}$
  (\code{clock\_edges.py}:71--72).
\item \textbf{H waveform peak-to-peak}: $\max\bar w-\min\bar w$.
\item \textbf{CDS offset}: on the median skip cycle $\bar c$, $\med\bar c[s_0+5..s_0+144]-\med\bar c[5..144]$: the CDS of an
  empty pixel, i.e.\ the summing-well feedthrough.
\item \textbf{Skip pulse}: the largest $|\bar c-\text{baseline}|$ outside the two windows.
\end{itemize}

\subsection{From traces to one feature row}
Every trace of a packet and amplifier is processed separately; each output column is then the \textbf{median over that
packet's traces} (3, sometimes 4), column by column (\code{physics\_features.py}:66--79; same in \code{gain\_features.py}).
So $\sigma_{500}$ of a packet is the median of the per-trace $\sigma_{500}$, not recomputed from median levels.
No drift correction, no standardisation at this stage; traces that fail are skipped.

%=====================================================================================================================
\section{Extremely randomised trees}\label{sec:ert}

Three of the main models are forests of extremely randomised trees (ERT) \cite{geurts2006}, as implemented in scikit-learn
\cite{pedregosa2011}. The method is written down once here.

\paragraph{One tree.} A regression tree partitions input space into boxes (leaves) by axis-aligned cuts and predicts the
mean training label of the box \cite{breiman1984}. It is grown from the root by recursively splitting a node holding the
training rows $\mathcal S$:
\begin{enumerate}
\item draw $m$ features at random without replacement (\code{max\_features});
\item for each drawn feature $j$, draw \emph{one} cut point uniformly, $t_j\sim\mathcal U\big[\min_{\mathcal S}x_j,\max_{\mathcal S}x_j\big]$;
\item keep the candidate split $(j,t_j)$ with the largest impurity decrease
\begin{equation}
\Delta(j,t)=|\mathcal S|\,\mathrm{Var}(\mathcal S)-|\mathcal S_L|\,\mathrm{Var}(\mathcal S_L)-|\mathcal S_R|\,\mathrm{Var}(\mathcal S_R),
\qquad \mathcal S_L=\{i\in\mathcal S:x_{ij}\le t\};
\end{equation}
\item stop when a child would hold fewer than \code{min\_samples\_leaf} rows (or the node is pure).
\end{enumerate}
The difference from a random forest \cite{breiman2001} is step 2 (random cut instead of the best cut) and that each tree
sees \emph{all} training rows (\code{bootstrap=False}). The random cut makes the trees more different from each other,
which lowers the variance of their average.

\paragraph{The forest.} $B$ trees $T_b$ grown with different random draws; the prediction is the average,
\begin{equation}
\hat y(x)=\frac1B\sum_{b=1}^{B}T_b(x).
\end{equation}

\paragraph{Classification.} For a label $u\in\{0,1\}$ the impurity is the Gini index $\mathrm{Gini}(\mathcal S)=2\hat p(1-\hat p)$ with
$\hat p$ the (weighted) fraction of $u=1$. With \code{class\_weight="balanced"} each row of class $c$ carries the weight
$n/(2n_c)$, so both classes weigh the same in total \cite{pedregosa2011}. The forest's probability is the average over trees of
the weighted class fraction in the leaf:
\begin{equation}
P(u=1\mid x)=\frac1B\sum_b \frac{\sum_{i\in\ell_b(x)}\omega_i u_i}{\sum_{i\in\ell_b(x)}\omega_i},
\end{equation}
with $\ell_b(x)$ the leaf of tree $b$ containing $x$. It equals the fraction of trees ``voting usable'' only when every
leaf is pure.

%=====================================================================================================================
\section{The noise model: trace physics + voltages, 8.6\,\% on a new amplifier}\label{sec:noise}

Code: \code{amp\_phys/test1\_physics\_model.py} (``test 1''), results \code{amp\_phys/results/test1/test1.md}.

\subsection{Task and split}
Predict $y_\sigma=\ln\sigma_e$ of an image on an amplifier \emph{never seen in training}, from the voltages and that packet's
traces, before the image is taken.
\begin{itemize}
\item Rows: board 101, channels 0, 1, 3, sessions G and H, image fit usable (\code{label.ok}), CDS window not
  fractional, $S_w>0$ and $\sigma_{500}$ finite (lines 32--49): 302 rows, 235 in G and 67 in H.
\item For each held-out amplifier $a\in\{0,1,3\}$: \textbf{train} on session G of the other two amplifiers, \textbf{test} on
  session H of amplifier $a$ (lines 112--120). Every test row is thus a new amplifier \emph{and} a later session.
\item No hyperparameter is tuned inside the test.
\end{itemize}

\subsection{Inputs}
\begin{itemize}
\item \textbf{Trace physics values} $\phi$ (9 numbers, all logged; \code{phys\_cols}, lines 56--61):
\begin{equation}
\phi=\big(\ln\sigma_{500},\ \ln S_w,\ \ln K,\ \ln\sigma_1(\tau_1),\dots,\ln\sigma_1(\tau_6)\big).
\end{equation}
\item \textbf{Voltages} $\theta$: the 31 \code{theta.*} columns that vary (s.d.\ $>10^{-6}$), raw volts, NaN$\to0$, not scaled
  (trees do not need scaling). Many are tied in groups, so they are not independent.
\item \textbf{One derived feature}, the physics prediction of $\ln\sigma_e$ without its offset:
  $\psi=\ln\sigma_{500}-\widehat{\ln g}(\theta)$, with $\widehat{\ln g}$ from the gain sub-forest below.
\end{itemize}
Total: $9+31+1=41$ inputs.

\subsection{The gain sub-forest}
The trace gives $\sigma_{500}$ in ADU; the label is in electrons, so a gain is needed. It is predicted from the voltages:
an ERT $\widehat{\ln g}(\theta)$ trained on the same G rows with target $\ln g$ (lines 72--73). On the training rows the
derived feature uses the in-sample prediction, on the test rows the out-of-sample one.

\subsection{The model}
\begin{equation}
\widehat{\ln\sigma_e}=\frac1{400}\sum_{b=1}^{400}T_b\big(\phi,\theta,\psi\big),
\end{equation}
an ERT regressor (Sec.~\ref{sec:ert}) with exactly these settings (line 53):
\begin{center}\code{ExtraTreesRegressor(n\_estimators=400, min\_samples\_leaf=3, max\_features=0.5, random\_state=0)}\end{center}
i.e.\ $B=400$ trees, $m=\lfloor0.5\cdot41\rfloor=20$ candidate features per split, at least 3 rows per leaf, squared-error
impurity, no depth limit, all rows to every tree. \textbf{Training} is the tree growing of Sec.~\ref{sec:ert} on the
$\approx155$ G rows of the two training amplifiers; there is no gradient descent and no validation-based early stopping.

\subsection{The three comparison models}
\begin{enumerate}
\item \textbf{Voltages only}: the same ERT on $\theta$ alone.
\item \textbf{Physics formula, gain from voltages} (``no ML'' on the noise side; lines 71--75):
\begin{equation}
\widehat{\ln\sigma_e}=\ln\sigma_{500}-\widehat{\ln g}(\theta)+c,\qquad
c=\med_{\rm train}\big[y_\sigma-\ln\sigma_{500}+\widehat{\ln g}(\theta)\big],
\end{equation}
i.e.\ $\hat\sigma_e=e^c\,\sigma_{500}/\hat g$. One number, $c$, is learned; in the three folds $e^c=0.995,\ 0.955,\ 0.944$.
\item \textbf{Trace physics values only}: the ERT on $\phi$ alone (9 inputs).
\end{enumerate}

\subsection{Score}
With $e_i=\widehat{\ln\sigma_{e,i}}-\ln\sigma_{e,i}$ pooled over the three held-out amplifiers (67 H rows),
\begin{equation}
\text{typical error}=100\sqrt{\tfrac1n\textstyle\sum_i e_i^2}\ \%\qquad\text{(RMS of the log error; lines 85--100).}
\end{equation}

\begin{center}\small
\begin{tabular}{lcccc}\toprule
model (new amplifier, sealed H) & typical error & median $|e|$ & change $\rho$ & change size error\\\midrule
voltages only & 18.2\,\% & 14.7\,\% & 0.65 & 12.3\,\%\\
physics formula, gain from voltages & 11.9\,\% & 7.4\,\% & 0.63 & 10.3\,\%\\
trace physics values only & 9.6\,\% & 5.5\,\% & 0.54 & 10.0\,\%\\
\textbf{trace physics values + voltages} & \textbf{8.6\,\%} & 5.7\,\% & 0.69 & 9.6\,\%\\\midrule
same, amplifier seen in training & 6.5\,\% & 2.3\,\% & 0.81 & 8.2\,\%\\\bottomrule
\end{tabular}
\end{center}
``Change'' columns: for every pair of H rows of one amplifier, Spearman's $\rho$ \cite{spearman1904} between the predicted and true
$\Delta\ln\sigma_e$, and the RMS of their difference, averaged over the amplifiers. Per held-out amplifier the best model
gives 8.6, 9.5 and 7.5\,\% (ch0, ch1, ch3).

\paragraph{Why it works.} Most of the variation of $\sigma_e$ between amplifiers is the amplifier's own noise, which the
voltages cannot know but the trace measures directly ($\sigma_{500}$ is a physical prediction of the same quantity). The
forest learns the remaining correction (gain, and what the 0.2\,s trace misses) from the voltages and the shape of
$\sigma_1(\tau)$.

\caveat{\textbf{Status of the 8.6\,\%.} (i) Test 1 was written on 17 November, after session H had been opened for the
scorecard (its docstring quotes the H scorecard). H is never used for fitting, but the design (inputs, ERT, its settings)
was chosen with H results known, so H is less strictly sealed for this model than for the scorecard's frozen models.
(ii) No uncertainty is computed in the repo. A packet bootstrap \cite{efron1979} over the 23 H packets, done for this note and
not stored in the repo, gives a 95\,\% interval of 6.2--11.2\,\% for the best model and 7.4--12.0\,\% for physics values only:
the two overlap heavily. Changing the forest seed moves the best model between 8.4 and 8.7\,\%. (iii) One board, one night,
three amplifiers. (iv) The gain feature uses in-sample predictions on the training rows, which are slightly better than the
test-row predictions.}

\caveat{\textbf{The forest inside the search is a different forest.} The trace-driven ranking and the hybrid
(Sec.~\ref{sec:strategies}) use an ERT with the same settings but on different inputs: the 18 standardised knobs and the 8
\code{PHYS\_COLS} of \code{fleet\_tests.py}:194--195 ($\sigma_{500}$, $\sigma_{500}^{\rm white}$, $S_w$, $K$, reset jump,
H peak-to-peak, skip pulse, CDS offset), with no $\sigma_1(\tau)$ and no gain sub-forest. Its error on a new amplifier has not been
scored separately; the 8.6\,\% belongs to the test-1 model.}

%=====================================================================================================================
\section{The usability classifier and the gate}\label{sec:gate}

About 22\,\% of the v3 settings (82 of 381 images) gave an image that could not be used. The search must avoid spending
images on them.

\subsection{What ``usable'' means}
\begin{equation}
u=\one\big[\code{label.ok}\big]\cdot\one\big[\ln\sigma_e,\ln\lambda,\ln g\ \text{finite}\big]\cdot\one\big[\lambda>0.005\ \text{e}^-/\text{pixel}\big]
\end{equation}
(\code{fleet\_tests.py}:108--109). \code{label.ok} is false when the WADERs fit raised any flag
(\code{process\_packet.py}:534--565): not fitted (flooded or unresolved peaks); a value $\le0$, an undefined error, or an error
larger than half the value for $g$, $\sigma_e$ or $\lambda$; $\chi^2/\text{ndf}>3$. The $\lambda$ floor removes fits that
converged on a dark charge 400$\times$ below normal. On v3: 70 of the 82 unusable images have ``$\lambda$ not constrained'',
9 have $\chi^2/\text{ndf}>3$, 3 pass the flags but have $\lambda<0.005$.

\caveat{Earlier text described unusable settings as ``electron peaks merge, most columns turn hot''. That is the physical
picture; the coded criterion is the fit flags above.}

\subsection{Inputs}\label{sec:knobs}
\paragraph{Knobs $z$} (\code{fleet\_tests.knobs}, lines 151--175): every varied \code{theta.*}; columns that are identical or
correlated with $|r|>0.999$ merge into one knob; timing knobs on a $\log_2$ scale; each knob z-scored over all loaded rows,
$z_j=(\theta_j-\bar\theta_j)/s_j$. On v3 this gives 18 knobs: SW\_L, RG\_H, RG\_L, V1--V3 high/low, H1--H3 high/low, VSUB, VDD, VR.

\paragraph{Trace features $\chi$} (\code{physics\_matrix}, lines 197--212): the 8 \code{PHYS\_COLS} ($\sigma_{500}$,
$\sigma_{500}^{\rm white}$, $S_w$, $K$, reset jump, H peak-to-peak, skip pulse, CDS offset); $\ln|\cdot|$ of each except the CDS
offset (kept linear, it changes sign), standardised over all rows, missing $\to0$.

\subsection{The model and its training}
An ERT classifier (Sec.~\ref{sec:ert}) for $u$:
\begin{center}\code{ExtraTreesClassifier(n\_estimators=B, min\_samples\_leaf=2, class\_weight="balanced", random\_state=0)}\end{center}
with the default $m=\lfloor\sqrt{p}\rfloor$ features per split and Gini impurity. Where it is used:
\begin{center}\small
\begin{tabularx}{\linewidth}{@{}>{\raggedright\arraybackslash}p{3.3cm}c>{\raggedright\arraybackslash}X>{\raggedright\arraybackslash}X@{}}\toprule
use & $B$ & inputs & training rows\\\midrule
replay gate, \code{p\_usable} (\code{fleet\_tests.py}:440--450) & 150 & those of the BO variant ($z$; or $[z,\chi]$ in the hybrid) & fleet (other amplifiers, other block; usable \emph{and} unusable) + every image taken so far on the new amplifier; \textbf{refitted after every image}\\
trace-driven ranking (lines 538--543) & 300 & $[z,\chi]$ & fleet only, once\\
live \code{propose\_next.py} & 300 & screen: 15 group deltas; image: + 8 trace features & every row of every amplifier\\\bottomrule
\end{tabularx}
\end{center}
If fewer than two unusable examples are available, $P\equiv1$.

\subsection{The hard gate}
With $P(x)=P(u=1\mid x)$ and a threshold $\gamma=0.5$, a candidate is allowed only if $P(x)\ge\gamma$:
\begin{equation}
\alpha(x)=\begin{cases}\EI(x) & P(x)\ge\gamma\\ -\infty & \text{otherwise}\end{cases}
\qquad\text{and if no candidate passes, image }\arg\max_xP(x)
\label{eq:gate}
\end{equation}
(\code{fleet\_tests.py}:613--615). The \emph{soft} alternative multiplies, $\alpha=\EI\cdot P$ (line 617). This is the
probability-of-feasibility weighting of constrained Bayesian optimisation \cite{schonlau1998,gardner2014,gelbart2014}.
The hard gate exists because broken settings often have unusual traces, which a Gaussian process sees as uncertain and
therefore as promising (large EI); a product can still let such a candidate win.

\subsection{How well it separates}
No script in the repo produces the gate's quoted scores; the obvious protocol was rebuilt for this note (train on G, test once
on H, 300 trees, the test amplifier left out, pooled; H has 20 unusable and 67 usable rows):
\begin{center}\small
\begin{tabular}{lccc}\toprule
inputs & AUC \cite{hanley1982} & unusable caught ($P<0.5$) & usable wrongly blocked\\\midrule
voltages only ($z$) & 0.990 & 11 of 20 & 0 of 67\\
voltages + trace ($[z,\chi]$) & 0.985 & 18 of 20 & 0 of 67\\\bottomrule
\end{tabular}
\end{center}
This reproduces the quoted ``18/20 with the trace, about 10/20 from voltages'' and ``0 usable blocked''. The quoted AUC of
0.99 is the voltages-only model's; with the trace it is 0.985. Both rank almost perfectly; the trace moves more of the unusable
settings below the 0.5 threshold.

\paragraph{Soft against hard in the replay} (images to 2\,\%, wasted images in brackets; nominal start;
\code{amp\_phys/results/fleet\_bo/fleet\_tests\_v3\_gate\_bo.md}):
\begin{center}\small
\begin{tabular}{lccc}\toprule
fleet BO variant & $w=1$ & $w=\frac12$ & $w=0$\\\midrule
voltages, soft & 2.5 (0) & 3.2 (0) & 3.0 (0)\\
voltages, hard gate & 3.7 (0) & 2.5 (0) & 2.7 (0)\\
voltages + trace, soft & 2.2 (0) & 2.5 (0) & 5.0 (2.2)\\
voltages + trace, hard gate & 2.2 (0) & 2.5 (0) & 3.3 (0.3)\\\bottomrule
\end{tabular}
\end{center}
The hard gate matters when the GP sees the traces; with voltages only the soft form already wasted no image.

%=====================================================================================================================
\section{Gaussian-process regression}\label{sec:gp}

\subsection{Model}
A Gaussian process \cite{rasmussen2006} puts a prior on functions: any finite set of values $f(x_1),\dots,f(x_n)$ is jointly
normal with mean 0 and covariance $K_{ij}=k(x_i,x_j)$. Observations are $y_i=\mu+f(x_i)+\varepsilon_i$,
$\varepsilon_i\sim\mathcal N(0,n^2)$. The kernel here is a scaled Matérn-5/2 with one length scale per input
(automatic relevance determination) \cite{rasmussen2006,stein1999}:
\begin{equation}
k(x,x')=s^2M(r),\qquad M(r)=\Big(1+\sqrt5\,r+\tfrac53r^2\Big)e^{-\sqrt5 r},\qquad r^2=\sum_{j}\frac{(x_j-x'_j)^2}{\ell_j^2}.
\end{equation}
$s^2$ is the variance of the response, $\ell_j$ how far knob $j$ must move to change it ($\ell_j\to\infty$ switches the knob
off), $n^2$ the image-to-image scatter at a fixed setting. Matérn-5/2 functions are twice differentiable: smooth, but less
rigid than a squared exponential.

\subsection{Prediction}
Given training inputs $X$ and labels $y$, with $K=k(X,X)$, $k_*=k(X,x_*)$ and $K_y=K+n^2I$, the posterior at a new
setting $x_*$ is normal with
\begin{equation}
m(x_*)=\mu+k_*^\top K_y^{-1}(y-\mu\one),\qquad
v(x_*)=k(x_*,x_*)-k_*^\top K_y^{-1}k_* .
\label{eq:gppost}
\end{equation}
$m$ is a weighted sum of the measured values, weighted by similarity; $v$ starts at $s^2$ and shrinks near measured settings.
The code solves with a Cholesky factor $K_y=LL^\top$ (plus $10^{-9}I$ for stability) and returns the latent standard deviation
$s(x_*)=\sqrt{v}$, without $n^2$.

\subsection{Training}
The hyperparameters $\eta=(s^2,\ell_1,\dots,\ell_p,n^2)$ are chosen by maximising the log marginal likelihood (type-II
maximum likelihood) \cite{rasmussen2006}:
\begin{equation}
\ln p(y\mid X,\eta)=-\tfrac12(y-\mu)^\top K_y^{-1}(y-\mu)-\tfrac12\ln|K_y|-\tfrac n2\ln2\pi,
\qquad \tfrac12\ln|K_y|=\textstyle\sum_i\ln L_{ii}.
\label{eq:lml}
\end{equation}
The first term rewards fitting the data, the second penalises a flexible (large-volume) prior; together they choose the
length scales without a validation set. Training fits about 20 numbers, not weights.

%=====================================================================================================================
\section{The fleet model: a multi-task Gaussian process}\label{sec:fleet}

Code: class \code{MultiTaskGP}, \code{lvl1/fleet\_tests.py}:252--292.

\subsection{Objective}\label{sec:objective}
\begin{equation}
J=w\,\ln\sigma_e+(1-w)\ln\lambda,\qquad w\in\{1\ (\text{noise only}),\ \tfrac12\ (\text{noise + dark charge}),\ 0\ (\text{dark charge only})\}.
\end{equation}
Lower is better. Only usable images enter the GP. In the replay the labels are optionally drift-corrected (\code{--drift},
lines 136--148): for each amplifier,
\begin{equation}
\tilde y_i=y_i-\mathrm{interp}_t\big(\text{the amplifier's nominal packets except }i\big)(t_i)+\overline{y_{\rm nom}},
\end{equation}
which removes the drift during the night and keeps the level. This uses the new amplifier's own nominal packets, so the
replay also reports runs without it.

\subsection{Model}
For amplifier $a$ at setting $x$ (knobs $z$ of Sec.~\ref{sec:knobs}),
\begin{equation}
J_a(x)=\mu+f(x)+c_a+g_a(x)+\varepsilon,
\qquad
\begin{aligned}
f&\sim\mathcal{GP}\big(0,\ \rho\,s^2M\big)&&\text{shared shape (the fleet)}\\
g_a&\sim\mathcal{GP}\big(0,\ (1-\rho)\,s^2M\big)&&\text{amplifier's own shape, independent per }a\\
c_a&\sim\mathcal N(0,c^2)&&\text{amplifier's own level}\\
\varepsilon&\sim\mathcal N(0,n^2)&&\text{image-to-image scatter}
\end{aligned}
\end{equation}
all independent. Its covariance between two observations is
\begin{equation}
\mathrm{Cov}\big(J_a(x),J_{a'}(x')\big)=s^2M(x,x')\big[\rho+(1-\rho)\,\one[a=a']\big]+c^2\,\one[a=a']+n^2\,\one[\text{same image}].
\label{eq:mtk}
\end{equation}
This is an intrinsic-coregionalisation multi-task GP \cite{bonilla2007} with task correlation $\rho$ for the shape plus a
per-task offset; its use for transferring an optimisation between related tasks follows \cite{swersky2013}. The two new
numbers mean:
\begin{itemize}
\item $\rho\in[0,1]$: how alike two amplifiers' responses are in shape. $\rho=1$: one shared response plus an offset.
  $\rho=0$: unrelated.
\item $c^2$: how much their overall levels differ.
\end{itemize}
A constant offset $c_a$ does not move the minimum, so if the shapes agree the fleet already ranks the settings of an
amplifier that has never been imaged; the first image then fixes its level.

\subsection{Training: two stages}
The fleet model has $p+5$ hyperparameters: $s^2$, the $p$ length scales $\ell_1,\dots,\ell_p$, $n^2$, $\rho$ and $c^2$ ($p=18$ knobs on v3).
They are fitted on the fleet's usable rows only,
\[
\mathcal D=\{(x_i,\,a_i,\,y_i)\}_{i=1}^{n},\qquad x_i\in\mathbb R^p\ \text{(standardised knobs)},\quad a_i\ \text{(amplifier)},\quad y_i=J_i,
\]
in two steps: first the shape of the response ($s^2,\ell,n^2$), then how that variance is split between amplifiers ($\rho,c^2$).

\paragraph{Stage 1: one GP for the whole fleet} (\code{fit\_shared\_gp}, \code{fleet\_tests.py}:235--244).
The amplifier labels are ignored, i.e.\ every row is treated as coming from one amplifier (the special case $\rho=1$, $c^2=0$ of
eq.~\eqref{eq:mtk}). With the centred labels $\tilde y_i=y_i-\bar y$, $\bar y=\frac1n\sum_iy_i$, the model is
\begin{equation}
\tilde{\mathbf y}\sim\mathcal N\big(\mathbf 0,\ K_1(\eta)\big),\qquad
\big[K_1(\eta)\big]_{ij}=s^2\,M\!\left(\sqrt{\textstyle\sum_{k=1}^{p}(x_{ik}-x_{jk})^2/\ell_k^2}\right)+n^2\,\delta_{ij},
\qquad \eta=(s^2,\ell_1,\dots,\ell_p,n^2).
\end{equation}
The hyperparameters maximise the log marginal likelihood
\begin{equation}
\hat\eta=\arg\max_{\eta\in\mathcal B}\ \mathcal L_1(\eta),\qquad
\mathcal L_1(\eta)=-\tfrac12\,\tilde{\mathbf y}^\top K_1(\eta)^{-1}\tilde{\mathbf y}-\tfrac12\ln\big|K_1(\eta)\big|-\tfrac n2\ln2\pi ,
\label{eq:stage1}
\end{equation}
over the box
\begin{equation}
\mathcal B:\quad 10^{-5}\le s^2\le10^{2},\qquad 0.05\le\ell_k\le10^{3},\qquad 10^{-6}\le n^2\le10 .
\end{equation}
\begin{itemize}
\item \textbf{Parametrisation.} scikit-learn optimises $\vartheta=\ln\eta$ (so the bounds are on $\ln s^2,\ln\ell_k,\ln n^2$)
  \cite{pedregosa2011}.
\item \textbf{Gradient.} With $K_1=LL^\top$ and $\boldsymbol\alpha=K_1^{-1}\tilde{\mathbf y}$,
  $\dfrac{\partial\mathcal L_1}{\partial\vartheta_m}=\tfrac12\operatorname{tr}\!\Big[\big(\boldsymbol\alpha\boldsymbol\alpha^\top-K_1^{-1}\big)\dfrac{\partial K_1}{\partial\vartheta_m}\Big]$
  \cite[eq.~5.9]{rasmussen2006}.
\item \textbf{Optimiser.} L-BFGS-B \cite{byrd1995} from two starting points, keeping the better optimum:
  \begin{itemize}
  \item the default $s^2=\mathrm{var}(y)$, $\ell_k=2$ for every knob, $n^2=0.1\,\mathrm{var}(y)$;
  \item one random point, $\vartheta\sim\mathcal U(\ln\mathcal B)$ (\code{n\_restarts\_optimizer=1}).
  \end{itemize}
\item \textbf{No priors.} $\mathcal L_1$ is maximised as it is; only the box limits the values. A length scale that runs to the upper
  bound $10^3$ (i.e.\ $10^3$ standard deviations of that knob) switches the knob off.
\end{itemize}
Output: $\hat s^2,\hat\ell,\hat n^2$.

\paragraph{Stage 2: split the variance between amplifiers} (\code{MultiTaskGP.\_\_init\_\_} and \code{\_loglik},
\code{fleet\_tests.py}:258--284). Now the amplifier labels are used. $\hat s^2,\hat\ell,\hat n^2$ are held fixed, the mean is fixed at
$\mu=\bar y$, and for a candidate pair $(\rho,c^2)$ the covariance of all fleet rows is eq.~\eqref{eq:mtk}:
\begin{equation}
\big[K_2(\rho,c^2)\big]_{ij}=\hat s^2\,\hat M_{ij}\big[\rho+(1-\rho)\,\one[a_i=a_j]\big]+c^2\,\one[a_i=a_j]+\hat n^2\,\delta_{ij},
\qquad \hat M_{ij}=M(x_i,x_j;\hat\ell).
\end{equation}
Its log marginal likelihood, computed with the Cholesky factor $K_2=L_2L_2^\top$ (the constant $-\frac n2\ln2\pi$ is the same for every
pair and is dropped), is
\begin{equation}
\mathcal L_2(\rho,c^2)=-\tfrac12\,\mathbf r^\top K_2^{-1}\mathbf r-\textstyle\sum_{i=1}^{n}\ln\big[L_2\big]_{ii},
\qquad \mathbf r=\mathbf y-\mu\one ,
\label{eq:stage2}
\end{equation}
since $\tfrac12\ln|K_2|=\sum_i\ln[L_2]_{ii}$; a pair whose $K_2$ is not positive definite gets $\mathcal L_2=-\infty$. The pair is
chosen by exhaustive search over a $6\times6$ grid, with $c^2$ measured in units of $\hat s^2$:
\begin{equation}
(\hat\rho,\hat c^2)=\arg\max_{\rho\in R,\ \kappa\in C}\ \mathcal L_2(\rho,\ \kappa\hat s^2),\qquad
R=\{0,\,0.5,\,0.75,\,0.9,\,0.97,\,1\},\quad C=\{0,\,0.1,\,0.3,\,1,\,3,\,10\}.
\end{equation}

\paragraph{What the two numbers mean once fitted.} At the same setting, two different amplifiers have
\begin{equation}
\mathrm{Var}\,J_a(x)=\hat s^2+\hat c^2+\hat n^2,\qquad
\mathrm{Cov}\big(J_a(x),J_b(x)\big)=\hat\rho\,\hat s^2\quad(a\ne b),\qquad
\mathrm{corr}=\frac{\hat\rho\,\hat s^2}{\hat s^2+\hat c^2+\hat n^2},
\end{equation}
so $\hat\rho$ is the fraction of the shape variance $\hat s^2$ that the amplifiers share, and $\hat c^2$ adds level differences on
top. For example, the noise objective of the table below ($\hat s^2=0.0145$, $\hat n^2=0.0066$, $\hat\rho=0.75$, $\hat c^2=0.0044$) gives a
correlation between two amplifiers at the same setting of $0.75\cdot0.0145/0.0255=0.43$.

\paragraph{Why two stages.} Stage 1 is the ordinary GP fit of Sec.~\ref{sec:gp} on all fleet rows together, which is well determined
(about 100 rows); stage 2 then decides only two numbers on a coarse grid, which is cheap and cannot run away. The price is that
$s^2,\ell,n^2$ are not re-optimised together with $\rho$ and $c^2$. Stage 1 assumes all amplifiers identical; when stage 2 picks
$c^2>0$, the variance of one amplifier becomes $\hat s^2+\hat c^2$ instead of $\hat s^2$ being reduced to make room, and when it picks
$\rho<1$, part of $\hat s^2$ moves to the amplifiers' own shapes with the length scales of the shared fit. The optimum is also only as
fine as the grid (e.g.\ $\rho$ between 0.5 and 0.75 cannot be represented).

The hyperparameters are fitted once on the fleet and \textbf{not refitted} during a search; each new image only enters the data.

\caveat{An earlier description says $\rho$ and $c$ are fitted ``with the same likelihood'', which suggests a joint fit. The
code does the two stages above. The same length scales $\hat\ell$ are used for the shared shape $f$ and for each amplifier's own
shape $g_a$.}

\subsection{Prediction for the new amplifier}
The training set is the fleet's usable rows plus the new amplifier's usable images so far, each with its amplifier label.
Equation~\eqref{eq:gppost} with the covariance~\eqref{eq:mtk} gives (lines 286--292)
\begin{equation}
m(x_*,a)=\mu+k_*^\top(K+n^2I)^{-1}(y-\mu\one),\qquad
v(x_*,a)=s^2+c^2-k_*^\top(K+n^2I)^{-1}k_*,
\end{equation}
where $k_*$ has entries $\rho s^2M(x_*,x_i)$ for fleet rows and $s^2M(x_*,x_i)+c^2$ for the new amplifier's own rows.
\begin{itemize}
\item \textbf{No image yet}: $m=\mu+\rho s^2M_*^\top(K+n^2I)^{-1}(y-\mu)$, the fleet's shape shrunk by $\rho$;
  the variance keeps the unknown level $c^2$ and the unknown own shape $(1-\rho)s^2$.
\item \textbf{After $k$ images}: the own rows are correlated with each other through $c^2$, so the first image mostly
  determines $c_a$; further images teach $g_a$ near where they were taken.
\end{itemize}

\caveat{An earlier description gives the predictive variance as $s^2-k_*^\top K^{-1}k_*$; in the multi-task case the code
correctly adds $c^2$.}

\paragraph{Fitted values} (v3, amplifier 101:3 left out, fleet from the other block; refitted for this note with the repo
code):
\begin{center}\small
\begin{tabular}{lcccc}\toprule
objective & $s^2$ & $n^2$ & $\rho$ & $c^2$\\\midrule
noise only ($w=1$) & 0.0145 & 0.0066 & 0.75 & 0.0044\\
noise + dark charge ($w=\frac12$) & 0.237 & 0.037 & 0.50 & 0\\
dark charge only ($w=0$) & 0.650 & 0.107 & 0.50 & 0\\\bottomrule
\end{tabular}
\end{center}
Most length scales go to the upper bound, i.e.\ those knobs are switched off; for dark charge only VSUB, H2\_H, VDD, RG\_H
(and VR at $w=0$) stay active. For the dark-charge objectives $c^2=0$ on this split: the levels do not differ beyond
what the shape term explains.

\subsection{Assumptions}
\begin{enumerate}
\item The response is a smooth (Matérn-5/2), stationary function of the standardised knobs, with one length scale per knob
  shared by all amplifiers.
\item Amplifiers differ by a constant level $c_a$ and an independent own shape $g_a$ with the same smoothness; $\rho$ and $c^2$
  are the same for every pair of amplifiers.
\item The scatter is Gaussian in log space and the same for every amplifier and setting (homoscedastic), $n^2$.
\item The fleet and the new amplifier are exchangeable: the new amplifier is ``another draw'' from the same population.
  On v3 the fleet is two sibling amplifiers of one board and one night, which favours the test.
\item Unusable settings carry no information about $J$ (they are excluded, and handled by the classifier).
\item The hyperparameters learned on the fleet hold for the new amplifier.
\end{enumerate}

%=====================================================================================================================
\section{Bayesian optimisation}\label{sec:bo}

\subsection{Expected improvement}
Let $b$ be the best (lowest) usable $J$ imaged so far on the new amplifier. Imaging $x$ improves on it by
$I(x)=\max\big(0,b-J(x)\big)$. Under the GP, $J(x)\sim\mathcal N(m,s^2)$, so with $z=(b-m)/s$
\begin{equation}
\EI(x)=\mathbb E[I(x)]=\int_{-\infty}^{b}(b-j)\,\frac1s\varphi\!\Big(\frac{j-m}{s}\Big)\dd j
=(b-m)\,\Phi(z)+s\,\varphi(z)
\label{eq:ei}
\end{equation}
\cite{jones1998} (\code{fleet\_tests.py}:435--437). The first term is large where the guess is good (exploit), the second
where the model is unsure (explore). No exploration offset $\xi$ is used. If no usable image exists yet, $b=\min_x m(x)$ over
the candidates. $s$ is the latent standard deviation (without $n^2$).

\subsection{The replay: what is searched and what is known}\label{sec:replay}
The v3 night is replayed as if the voltages of one amplifier were being tuned live (\code{fleet\_tests.test\_bo}, lines 453--641).
Every image $i$ of the night has a setting $x_i$ (the 18 standardised knobs, Sec.~\ref{sec:knobs}), an amplifier $a_i$, a time
block $\beta_i$, a usability label $u_i\in\{0,1\}$ (Sec.~\ref{sec:gate}) and, if usable, an objective value $J_i$
(Sec.~\ref{sec:objective}). Packets in time order are cut into two contiguous blocks, $\beta\in\{1,2\}$ (lines 215--221).

One search is defined by a \emph{new} amplifier $a^\star$ and a block $\beta^\star$ (lines 467--469):
\begin{align}
\text{candidate pool}\quad &\mathcal P=\{i:\ a_i=a^\star,\ \beta_i=\beta^\star\}, && |\mathcal P|=63\text{--}64\ \text{(unusable ones included)},\\
\text{fleet}\quad &\mathcal F=\{i:\ a_i\ne a^\star,\ \beta_i\ne\beta^\star\}, && \text{usable part }\mathcal F_u=\{i\in\mathcal F:u_i=1\}.
\end{align}
The fleet is the other two amplifiers in the \emph{other} half of the night, so the new amplifier never shares a moment (drift,
setup state) with its training data. $3$ amplifiers $\times$ $2$ blocks $=6$ searches.

Before the search the code knows: every $x_i$ of the pool (the settings that could be chosen), and $(x_i,a_i,u_i,J_i)$ of the
fleet. It does \emph{not} know $u_i$ or $J_i$ of any pool image until that image is ``taken''; taking an image means looking up
its stored $u_i$ and $J_i$.

\subsection{The loop, step by step}\label{sec:loop}
The state of the search after $k$ images is the ordered list of images taken, $\mathcal R_k=(r_1,\dots,r_k)\subset\mathcal P$, and
its usable part $\mathcal R_k^u=\{r\in\mathcal R_k:u_r=1\}$.

\paragraph{Step 0: start} (lines 476--477). If the nominal setting is known, the first image is the earliest usable nominal image of
the pool, $\mathcal R_1=(r_1)$; otherwise the search starts empty, $\mathcal R_0=()$.

\paragraph{Step 1: fit the fleet model once} (line 595). Fit $\hat s^2,\hat\ell,\hat n^2,\hat\rho,\hat c^2$ on $\mathcal F_u$ by the two
stages of Sec.~\ref{sec:fleet}, and fix the mean $\mu=\frac1{|\mathcal F_u|}\sum_{i\in\mathcal F_u}J_i$. These are \emph{not}
refitted during the search.

\paragraph{Step 2: repeat while $k<|\mathcal P|$} (lines 601--618). With the untested candidates
$\mathcal C_k=\mathcal P\setminus\mathcal R_k$:
\begin{enumerate}[label=(\alph*),leftmargin=2em]
\item \textbf{Training data of the GP.} The fleet's usable images plus the new amplifier's usable images so far, each with its
  amplifier label:
  \begin{equation}
  \mathcal D_k=\{(x_i,a_i,J_i):i\in\mathcal F_u\}\ \cup\ \{(x_r,a^\star,J_r):r\in\mathcal R^u_k\}.
  \end{equation}
  Unusable images of the new amplifier carry no $J$ and are left out.
\item \textbf{Prediction for every candidate} $c\in\mathcal C_k$, by eq.~\eqref{eq:gppost} with the fleet covariance~\eqref{eq:mtk}
  and amplifier label $a^\star$:
  \begin{equation}
  m_k(c)=\mu+\mathbf k_c^\top\big(K+\hat n^2I\big)^{-1}(\mathbf J-\mu\one),\qquad
  s_k(c)=\sqrt{\hat s^2+\hat c^2-\mathbf k_c^\top\big(K+\hat n^2I\big)^{-1}\mathbf k_c},
  \end{equation}
  where $K$ is the covariance~\eqref{eq:mtk} of the points of $\mathcal D_k$ without the $\hat n^2$ term, $\mathbf J$ their
  objective values, and $\mathbf k_c$ has entries $\hat\rho\hat s^2M(x_c,x_i)$ for fleet points and $\hat s^2M(x_c,x_r)+\hat c^2$ for the
  new amplifier's own points.
\item \textbf{Incumbent} (line 609): the best usable value found so far, or, before any usable image, the best prediction,
  \begin{equation}
  b_k=\begin{cases}\min_{r\in\mathcal R_k^u}J_r & \mathcal R_k^u\neq\emptyset\\[2pt] \min_{c\in\mathcal C_k}m_k(c) & \text{otherwise.}\end{cases}
  \end{equation}
\item \textbf{Expected improvement} of each candidate, eq.~\eqref{eq:ei}:
  \begin{equation}
  \EI_k(c)=\big(b_k-m_k(c)\big)\,\Phi\big(z_k(c)\big)+s_k(c)\,\varphi\big(z_k(c)\big),\qquad z_k(c)=\frac{b_k-m_k(c)}{s_k(c)} .
  \end{equation}
\item \textbf{Usability probability} (\code{p\_usable}, lines 440--450). An ERT classifier (150 trees, Sec.~\ref{sec:gate}) is
  refitted from scratch on all fleet images, usable or not, and every image taken so far:
  \begin{equation}
  \mathcal U_k=\{(x_i,u_i):i\in\mathcal F\}\ \cup\ \{(x_r,u_r):r\in\mathcal R_k\},\qquad P_k(c)=P(u=1\mid x_c)\ \text{for }c\in\mathcal C_k ;
  \end{equation}
  if $\mathcal U_k$ holds fewer than two unusable images, $P_k\equiv1$. Unlike the GP, the classifier learns from the new
  amplifier's \emph{failed} images too.
\item \textbf{Gate and choice} (lines 613--615, 618):
  \begin{equation}
  \alpha_k(c)=\begin{cases}\EI_k(c)&P_k(c)\ge0.5\\-\infty&P_k(c)<0.5\end{cases},\qquad
  r_{k+1}=\begin{cases}\arg\max_{c\in\mathcal C_k}\alpha_k(c) & \text{some }P_k(c)\ge0.5\\[2pt]
  \arg\max_{c\in\mathcal C_k}P_k(c) & \text{none passes.}\end{cases}
  \end{equation}
\item \textbf{Take the image}: append $r_{k+1}$ to $\mathcal R$, reveal $u_{r_{k+1}}$ and, if usable, $J_{r_{k+1}}$; set $k\leftarrow k+1$.
\end{enumerate}
The loop runs through the whole pool (there is no stopping rule in the replay), so the full sequence $r_1,r_2,\dots$ is recorded and
scored afterwards. The soft variant replaces (f) by $r_{k+1}=\arg\max_c\EI_k(c)\,P_k(c)$.

\subsection{Scoring a search}
With $J^\star=\min_{i\in\mathcal P,\,u_i=1}J_i$ the best usable value in the pool (known to the scorer, not to the search), the gap after
$k$ images is
\begin{equation}
\Delta_k=100\Big(\min_{r\in\mathcal R^u_k}J_r-J^\star\Big)\ \%,\qquad \Delta_k=100\ \text{if }\mathcal R^u_k=\emptyset
\end{equation}
(\code{regret\_curve}, lines 480--483; an unusable image counts as $J=+\infty$, i.e.\ it never improves the best). Then
\begin{equation}
k_{2\%}=\min\{k:\Delta_k\le2\},\qquad \text{wasted}=\big|\{r\in\mathcal R_{k_{2\%}}:u_r=0\}\big|,
\end{equation}
counting the start image as image 1. The tables report the mean of $k_{2\%}$ and of the wasted count over the 6 searches.

\subsection{Strategies compared}\label{sec:strategies}
\begin{itemize}
\item \textbf{Person, hill-climb}: from the nominal (or a random start; mean over $\le50$ starts) image the untested setting
  nearest in standardised knobs to the best usable so far.
\item \textbf{Random order}: 200 random orders.
\item \textbf{Fleet ranking, no updates}: candidates in order of the zero-image $m$.
\item \textbf{Trace-driven} (lines 525--547): $\hat J=w\,\widehat{\ln\sigma_e}+(1-w)\widehat{\ln\lambda}$, with
  $\widehat{\ln\sigma_e}$ from an ERT (400 trees, leaf $\ge3$, $m=0.5p$) on $[z,\chi]$ trained on the fleet's usable rows, and
  $\widehat{\ln\lambda}$ from the zero-image fleet GP. Candidates with $P\ge0.5$ (300-tree classifier on $[z,\chi]$) are imaged in
  order of $\hat J$, the rest after them; images never change the order.
\item \textbf{Fleet BO + hard gate}: the loop above with $x=z$ and the gate on $z$.
\item \textbf{Hybrid} (lines 550--565): first image = the trace-driven top pick; then fleet BO on $z$ with the hard gate of the
  trace-aware classifier on $[z,\chi]$.
\item \textbf{BO without the fleet}: the same hyperparameters, but only the new amplifier's own data (after two usable random images).
\end{itemize}

\subsection{Results}
\begin{center}\small
\begin{tabular}{lccc}\toprule
images to 2\,\% (wasted) & noise only & noise + dark charge & dark charge only\\\midrule
\multicolumn{4}{l}{\textit{Nominal setting known}}\\
Person: hill-climb from nominal & \textbf{1.5} (0) & 26.8 (4.3) & 26.0 (4.5)\\
Random order & 5.0 (0.9) & 31.0 (4.7) & 38.0 (5.7)\\
Fleet ranking, no updates & 3.7 (0.2) & 9.2 (1.0) & 9.2 (1.0)\\
Trace-driven & 1.7 (0) & 8.3 (0) & 8.2 (0)\\
Fleet BO + hard gate & 3.7 (0) & \textbf{2.5} (0) & \textbf{2.7} (0)\\
Hybrid & 3.0 (0) & 2.8 (0) & \textbf{2.7} (0)\\
BO without the fleet & 5.8 (1.0) & 4.8 (0.7) & 4.8 (0.7)\\\midrule
\multicolumn{4}{l}{\textit{Nothing known about the new amplifier}}\\
Person: hill-climb from a random start & 11.5 (4.4) & 24.0 (3.1) & 26.8 (3.4)\\
Random order & 6.5 (1.5) & 31.0 (4.8) & 38.5 (5.7)\\
Fleet ranking, no updates & 4.3 (0.5) & 8.3 (1.0) & 8.3 (1.0)\\
Trace-driven & \textbf{1.3} (0) & 7.7 (0) & 7.3 (0)\\
Fleet BO + hard gate & 3.5 (0.2) & \textbf{1.7} (0) & \textbf{1.8} (0)\\
Hybrid & 2.8 (0) & 2.2 (0) & \textbf{1.8} (0)\\
BO without the fleet & 8.2 (2.2) & 3.2 (0.3) & 4.0 (0.5)\\\bottomrule
\end{tabular}
\end{center}
(\code{amp\_phys/results/fleet\_bo/fleet\_tests\_v3\_hybrid\_\{nominal,none\}\_w*\_bo.md}.) With 6 searches, differences below about
one image are not meaningful; ``BO without the fleet'' at $w=1$ is 5.8 in one file and 7.2 in another, only because the
random stream differs.

\paragraph{Reading.} Noise: the trace measures it, so ranking by the trace (or stepping out from a good nominal) wins.
Dark charge: the trace cannot see it, but the fleet knows the best dark charge is at low VSUB, far from nominal, so fleet BO
with the gate wins (1.7--2.7 images against 24--38 by hand). The hybrid is never worse than 3 images. With the gate no
image was wasted on a broken setting.

\subsection{The live procedure (\code{lvl1/propose\_next.py})}
In a live search the candidates are generated, not looked up. One round:
\begin{enumerate}
\item \textbf{Screen} (lines 173--214): draw 600 settings by scrambled Sobol sampling \cite{sobol1967,owen1998} over the 15 voltage
  groups, each within its allowed box on a 0.05\,V grid; redraw any that break the hardware windows and relational rules
  (\code{configs/limits.toml}, \code{src/ccd\_limits.py}), the 0.5\,V margins and 1.5\,V minimum swings, the ramp rule
  (VSUB in steps $\le5$\,V) or the reset-safety check. Score each by $\EI\cdot P$ with the hard gate. Keep up to 12 that are
  at least 0.75 apart in standardised space and write them as trace-only packets.
\item \textbf{Image} (lines 217--265): after the traces, re-score the screened settings with the trace-aware classifier, taking
  for each setting the \emph{lowest} $P$ over the board's amplifiers; image the best one.
\item \textbf{Status} (lines 268--288): ``left'' $=100\max_x\EI(x)\one[P(x)\ge\gamma]$ over fresh draws; stop if it is below
  \code{--stop} (default 1\,\%).
\end{enumerate}
Here the inputs are 15 group deltas from nominal (a swing group as high $-$ low minus its nominal), z-scored; the timing is
pinned to nominal; the GP is refitted on every usable row (fleet and this board) at each call; there is no drift correction.
For a whole board the objective is the mean over its amplifiers, $\bar m=\frac1A\sum_am_a$, with
$\bar s=\sqrt{\sum_as_a^2}/A$.

\caveat{(i) $\bar s$ treats the amplifiers' errors as independent and ignores their shared $f$, so it underestimates the
uncertainty of the board mean. (ii) Earlier text says the search stops ``when the expected gain is below the repeatability'';
the code compares with a fixed 1\,\% (the repeatabilities 1.7\,\% for $\sigma_e$ and 4\,\% for $\lambda$ are only printed).
(iii) The live code uses 15 knobs, the replay 18. (iv) The dry-run numbers quoted elsewhere (stop at 0.25\,\%; 8 screened settings
with $P=0.89$--$0.97$) are not stored in the repo.}

%=====================================================================================================================
\section{What is established and what is not}
\begin{itemize}
\item \textbf{Established on v3}: the trace halves the noise error on a new amplifier (18.2 $\to$ 8.6\,\%); the gate with the
  trace catches 18 of 20 unusable settings and blocks no usable one; fleet BO with the gate reaches 2\,\% of the best in about
  two images for dark-charge objectives.
\item \textbf{Limits}: one board, one night, three sibling amplifiers; the fleet and the new amplifier share the module and the
  night. On v2 + v3 together one objective took 16--18 images, because trace features shift between nights. The next test is a
  live run on a CCD the fleet has never seen.
\end{itemize}

%=====================================================================================================================
\appendix
\section{The other models}\label{app:other}

Two pieces of code produced the v3 numbers: the \textbf{scorecard} (\code{lvl1/scorecard\_v2.py}, with the model registry in
\code{lvl1/scorecard.py} and \code{lvl1/train\_lvl1.py}; results \code{processed\_v3n/lvl1\_results/scorecard\_\{G,H\}}) and the
\textbf{physics tests} (\code{amp\_phys/test1\_physics\_model.py}, \code{test2\_gain\_lambda.py}). Section~\ref{sec:noise} is test 1.

\subsection{Protocol of the scorecard}
\begin{itemize}
\item Rows: board 101 ch0/1/3, label usable, trace usable, window not fractional, $\lambda>0.005$.
\item Model choice on session G by cross-validation: 5 contiguous time blocks; for ``new amplifier'', leave one amplifier out
  $\times$ block. The top 3 per target and cell were \textbf{frozen} (16 November, 07:06,
  \code{lvl1/frozen\_choice\_v3.json}), together with the physics, settings and baseline models; then H was scored once.
\item Columns with $>20$\,\% missing or zero spread are dropped, gaps get the training median; each learner standardises its
  inputs on the training rows. \textbf{Trace models never see the voltages} (\code{no\_settings}).
\item Absolute error: RMS of the log error, as in Sec.~\ref{sec:noise}. Change models: size error $=100\,\mathrm{RMS}(\widehat{\Delta y}-\Delta y)$
  and Spearman $\rho$ over all scored pairs. Better/worse/same: a packet bootstrap with Poisson(1) weights per packet
  \cite{efron1979,hanley2006}, 1000 draws, on the difference of RMS to a reference model.
\end{itemize}

\subsection{Learners}
\paragraph{Gaussian process} (\code{src/ml/models.py}:114--159): $k=c\,M_{5/2}(\text{ARD})+\sigma_n^2\delta$,
$c\in[10^{-3},10^3]$, $\ell\in[10^{-2},10^3]$, $\sigma_n^2\in[10^{-6},10]$, \code{normalize\_y}, L-BFGS-B on eq.~\eqref{eq:lml}
for at most 150 iterations, no restarts, at most 400 training points (800 for change models).

\paragraph{Gradient-boosted trees} (``trees'' in the scorecard; \code{models.py}:80--111). An additive model built stagewise
\cite{friedman2001},
\begin{equation}
F_0=\bar y,\qquad F_t(x)=F_{t-1}(x)+\nu\,h_t(x),\qquad h_t=\arg\min_h\sum_i\big(r_i^{(t)}-h(x_i)\big)^2,\quad r_i^{(t)}=y_i-F_{t-1}(x_i),
\end{equation}
each $h_t$ a tree on histogram-binned inputs (scikit-learn \code{HistGradientBoostingRegressor} \cite{pedregosa2011,ke2017})
with at most 15 leaves and 10 rows per leaf; $\nu=0.05$, 300 stages. Two extra models with the pinball (quantile) loss
\cite{koenker1978} at 0.16 and 0.84 give the error bar $\hat s=|q_{84}-q_{16}|/2$.

\paragraph{Neural-network ensemble} (\code{models.py}:289--329). Linear($p$,64)--SiLU--Linear(64,64)--SiLU \cite{elfwing2018} with two
heads $(\mu,\ln\sigma^2)$, trained on the Gaussian negative log-likelihood $\tfrac12[\ln\sigma^2+(y-\mu)^2/\sigma^2]$
\cite{nix1994}; 5 networks \cite{lakshminarayanan2017}, AdamW \cite{loshchilov2019} (learning rate $3\times10^{-3}$, weight decay $10^{-3}$),
batch 64, up to 600 epochs, early stopping on 15\,\% of packets (patience 60), then refit on all rows. Prediction: mean of means;
variance = mean variance + variance of means.

\paragraph{1-D CNN on the trace curves} (\code{models.py}:332--396). Input: the trace curves (waveform shapes, Allan deviation,
CDS spectrum; own and reference trace) as channels, 256 points each. Conv1d($C$,16,7)--SiLU--MaxPool2 $\to$ Conv1d(16,32,7)--SiLU--MaxPool2
$\to$ Conv1d(32,32,5)--SiLU--AdaptiveAvgPool(8), concatenated with scalar features, Linear(256+$p$,64)--SiLU, heads $(\mu,\ln\sigma^2)$;
3 networks, learning rate $2\times10^{-3}$, batch 256, 400 epochs, patience 50.

\paragraph{Pretrained CNN} (\code{src/ml/deep.py}:68--229). The convolution stack plus Linear(256,16) is first trained without labels
for 40 epochs on: masked reconstruction (25\,\% of 32-sample blocks, or one whole curve), plus $0.3\times$ InfoNCE
\cite{oord2018} (temperature 0.2) between two augmented views ($\pm5$\,\% scale, noise 0.05, shift $\pm2$); then fine-tuned with
the labels (learning rate $5\times10^{-4}$).

\paragraph{Deep-kernel GP} (\code{deep.py}:232--350) \cite{wilson2016}: $k(i,j)=s^2M_{5/2}(|z_i-z_j|/\ell)+\sigma_n^2\delta_{ij}$ on
$z=$ (standardised frozen pretrained encoding, scalar features); Adam on the exact negative log marginal likelihood, 250 steps.

\paragraph{SAAS GP} (for $\lambda$ inside \code{trace: physics}; \code{src/ml/bogp.py}:73--77). A GP whose inverse squared length
scales have half-Cauchy priors that shrink most inputs off, $\tau\sim\mathcal{HC}(0.1)$, $\ell_j^{-2}\sim\mathcal{HC}(\tau)$
\cite{eriksson2021}; sampled by NUTS \cite{hoffman2014} (96 warm-up, 48 samples, thinning 8).

\subsection{Absolute models}
\begin{itemize}
\item \textbf{mean}, \textbf{amplifier mean}: the training mean (per amplifier).
\item \textbf{settings: gp / trees}: $\theta\to y$. \textbf{settings + amplifier}: $[\theta,\text{one-hot}(a)]\to y$ (seen amplifiers only).
\item \textbf{trace: gp / trees / nn}: own- and reference-trace features $\to y$ (feature tiers: gp the physics tier P, trees
  all, nn P+A).
\item \textbf{trace: physics} (\code{physics\_white}; \code{src/ml/structured.py}:655--704):
\begin{equation}
\widehat{\ln\sigma_e}=\tfrac12\ln\frac{w^2}{N}-\widehat{\ln g}(x)+\hat r(x),\qquad \hat s^2=s_g^2+s_r^2,
\end{equation}
$w$ the trace's fitted white noise per skip, $N$ the number of skips; $\widehat{\ln g}$ an ARD Matérn-5/2 GP on $\ln g$;
$\hat r$ a GP on the residual $r=\ln\sigma_e-(\tfrac12\ln w^2/N-\ln g)$, with rows of robust $|z|>4$ down-weighted by an
extra variance $(r-\med r)^2$; the labels' own variances on the diagonal. $\lambda$ from the SAAS GP. Inputs include the
\emph{packet state}, the mean over the packet's amplifiers of the reference-trace features minus each amplifier's average.
\item \textbf{trace: cnn / pretrained cnn / deep kernel}: as above, on the curves.
\item \textbf{Test 2} (\code{test2\_gain\_lambda.py}). Gain, reset formula: $\ln g=\ln|J_{\rm reset}|-\ln(RG_H-VR-V^*)+c$ by least squares
  ($V^*\approx9.3$\,V); gain, SW formula: $\ln|\text{CDS offset}|-\ln(\text{SW swing})+c$; ERT on the gain signals with and
  without $\theta$. $\lambda$: physics form $\lambda=e^{a+b\tilde V_{\rm SUB}}+e^{c+\mathbf d\cdot[\tilde H_{\rm sw},\tilde V_{\rm sw},\widetilde{SW}_{\rm sw},\widetilde{\rm OD}]}$
  (least squares in log space); ERT on dark and CIC separately; ERT on $\theta$ + swings + physics + gain signals (``settings + trace + setup'').
\end{itemize}

\subsection{Change models}
Pairs $A\to B$ of one amplifier, target $\Delta y=y_B-y_A$; random-pair labels drift-corrected against the nominal packets;
training uses both orders $(A,B)$ and $(B,A)$ with $\Delta\to-\Delta$. $F$ = trace features (logged where positive), $C$ = curves.
\begin{itemize}
\item \textbf{settings change: gp / trees}: $\Delta\theta\to\Delta y$.
\item \textbf{trace change: trees}: $F_B-F_A\to\Delta y$.
\item \textbf{combined trees} $[A,B,B-A]$: GBT $h$ on $[F_A,F_B,F_B-F_A]$, antisymmetrised,
  $\widehat\Delta=\tfrac12\big(h(A,B)-h(B,A)\big)$; \textbf{settings + trace combined}: the same with $[\theta_A,\theta_B,\Delta\theta]$ added.
\item \textbf{residual trees}: $b=\text{GBT}(\Delta\theta)$ out-of-fold, then an antisymmetrised GBT on $[\Delta\theta,F_A,F_B,\Delta F]$
  fitted to $\Delta y-b$.
\item \textbf{physics difference}: $\widehat{\Delta\ln\sigma_e}=\Delta\ln w-\widehat{\Delta\ln g}$, $\widehat{\Delta\ln g}=\text{GBT}(\Delta F)$;
  \textbf{+ correction}: a GBT on [base, $\Delta F$] fitted to the residual.
\item \textbf{shared-encoder net}: $\widehat\Delta=g(F_B)-g(F_A)$, $g$ = MLP($p$,32,32,3); loss Huber \cite{huber1964} on the
  scaled change plus a sign term $\overline{\omega\,\mathrm{softplus}(-\mathrm{sign}(\Delta)\,3\widehat\Delta)}$; 5 networks.
\item \textbf{difference curves PCA + trees}: $D=C_B-C_A$, 10 principal components \cite{pearson1901} fitted on $\{D,-D\}$, then GBT;
  \textbf{hand-picked bands + gp}: for each target the 6 curves whose standardised difference correlates most with $\Delta y$, averaged
  over 9 points around each peak, plus $\Delta\ln w$, into a GP.
\item \textbf{direction classifiers}: a gradient-boosted classifier for $\one[\Delta>0]$, weighted by $1-e^{-\Delta^2/2s^2}$
  ($s$ = label noise of a change), output $P-\frac12$ antisymmetrised; scored only on direction and $\rho$.
\item \textbf{absolute, differenced}: $\hat m(B)-\hat m(A)$.
\item \textbf{ensembles}: the mean of four members. Trace: trace change trees, combined trees, physics difference + correction,
  absolute trace physics differenced. Settings + trace: settings + trace combined trees, residual trees, physics difference +
  correction, absolute trace physics differenced (for $\lambda$ and $g$ the $\sigma_e$-only member is absent).
\end{itemize}

\subsection{Results on sealed session H}
\paragraph{Absolute, new amplifier} (typical error, \%):
\begin{center}\small
\begin{tabularx}{\linewidth}{@{}lXX@{}}\toprule
target & physics tests (G $\to$ H, ERT) & scorecard (frozen)\\\midrule
$\sigma_e$ & voltages 18.2 $\to$ physics formula 11.9 $\to$ physics values 9.6 $\to$ \textbf{physics + voltages 8.6} &
  trace trees 12.7, trace physics 13.5, trace nn 14.2, settings gp 18.8, settings trees 18.1\\
$\lambda$ & voltages 69.0 $\to$ \textbf{settings + trace + setup 48.8} (physics form 71.5, dark + CIC 83.2) &
  trace trees 49.5, trace physics 50.9, trace nn 51.1, settings gp 71.7, settings trees 66.4\\
$g$ & voltages 3.61, physics values 3.53, \textbf{physics + voltages 3.42} (reset formula 8.7, SW formula 69.4) &
  settings trees 3.6, settings gp 3.9, trace physics 4.3\\\bottomrule
\end{tabularx}
\end{center}
Not frozen (G cross-validation only, new amplifier, $\sigma_e$ / $\lambda$): cnn 20.4 / 59.9, pretrained cnn 18.7 / 63.2,
deep kernel 23.1 / 84.0, trace gp 17.9 / 66.3: the flexible models overfit the $\approx$155 training rows.

\paragraph{Absolute, amplifier seen in training} (scorecard H): $\sigma_e$: trace physics \textbf{7.6}, settings + amplifier trees 7.7,
trace trees 8.0 (test 1's physics + voltages forest: 6.5). $\lambda$: settings + amplifier trees \textbf{10.5}, settings + amplifier gp
14.6, trace physics 24.2. $g$: settings + amplifier gp \textbf{0.6}, settings + amplifier trees 0.9, trace physics 1.1.

\paragraph{Change, new amplifier, random pairs} (scorecard H; $\rho$ / size error \%):
\begin{center}\small
\begin{tabularx}{\linewidth}{@{}lccX@{}}\toprule
target & no change & settings change trees & best, then runners-up\\\midrule
$\sigma_e$ & -- / 14.0 & 0.64 / 12.9 & \textbf{ensemble settings + trace 0.73 / 11.2}; physics difference + correction 0.65 / 11.5; trace ensemble 0.63 / 11.5\\
$\lambda$ & -- / 63.5 & 0.89 / 38.9 & \textbf{settings + trace combined trees 0.94 / 29.9}; residual trees 0.91 / 33.5; ensemble 0.91 / 34.3\\
$g$ & -- / 3.4 & \textbf{0.87 / 1.8} & settings + trace ensemble 0.87 / 2.1; combined trees 0.86 / 1.9\\\bottomrule
\end{tabularx}
\end{center}
Not frozen ($\sigma_e$, G only): shared-encoder net $\rho=0.36$, hand-picked bands + gp 0.55.

\paragraph{Change, amplifier seen}: $\sigma_e$ ensemble settings + trace $\rho=0.82$ / 10.0; $\lambda$ settings + trace combined trees
0.98 / 15.3; $g$ settings + amplifier gp differenced 0.97 / 0.8.


\subsection{The extremely randomised trees on every input type}\label{app:ertinputs}
The forest of Sec.~\ref{sec:noise} (400 trees, $\ge3$ rows per leaf, half the features per split), trained on session G and
scored once on session H with every combination of input types (\code{amp\_phys/ert\_inputs.py}, results
\code{amp\_phys/results/ert\_inputs/}). The inputs:
\begin{itemize}
\item \textbf{voltages}: the 31 varied \code{theta.*} columns;
\item \textbf{physics values}: the 9 of Sec.~\ref{sec:noise}, $\ln\sigma_{500},\ln S_w,\ln K,\ln\sigma_1(\tau_{1..6})$;
\item \textbf{signal sizes}: the 4 of Sec.~\ref{sec:gainfeat}, $\ln|\cdot|$ of reset jump, H peak-to-peak and skip pulse, and the CDS offset;
\item \textbf{141 trace features}: the general features of each packet's own traces (\code{trace.*}), columns with $>20$\,\% missing
  or no spread dropped, gaps filled with the training median.
\end{itemize}
``New'': the amplifier is left out of training; ``seen'': all three are in training, with no amplifier identity as input.
Rows as in test 1 plus $\lambda>0.005$ (67 H rows for every target). Typical error (RMS of the log error, \%) with a 95\,\%
packet bootstrap interval (Poisson(1) weight per H packet, 1000 draws \cite{efron1979,hanley2006}); best per column in bold.
\begin{center}\footnotesize
\setlength{\tabcolsep}{3pt}
\begin{tabular}{@{}lcccccc@{}}\toprule
 & \multicolumn{2}{c}{$\sigma_e$} & \multicolumn{2}{c}{$\lambda$} & \multicolumn{2}{c}{gain}\\
\cmidrule(lr){2-3}\cmidrule(lr){4-5}\cmidrule(lr){6-7}
inputs & new & seen & new & seen & new & seen\\\midrule
voltages (31) & 18.4 {\scriptsize[16.1, 21.2]} & 13.6 {\scriptsize[10.9, 17.3]} & 66.8 {\scriptsize[58.1, 75.1]} & 46.9 {\scriptsize[41.1, 53.0]} & 3.6 {\scriptsize[3.4, 3.9]} & 2.5 {\scriptsize[2.4, 2.7]}\\
physics values (9) & 9.1 {\scriptsize[7.4, 11.7]} & 7.0 {\scriptsize[4.9, 9.9]} & 50.8 {\scriptsize[43.7, 59.9]} & 31.2 {\scriptsize[23.1, 41.4]} & 4.1 {\scriptsize[3.7, 4.6]} & 2.7 {\scriptsize[2.2, 3.2]}\\
signal sizes (4) & 14.8 {\scriptsize[12.3, 18.0]} & 9.7 {\scriptsize[6.9, 13.4]} & 52.1 {\scriptsize[42.9, 62.5]} & 39.4 {\scriptsize[28.3, 51.7]} & 3.9 {\scriptsize[3.6, 4.2]} & 2.5 {\scriptsize[2.0, 3.0]}\\
141 trace features & 16.2 {\scriptsize[14.8, 17.8]} & 7.0 {\scriptsize[4.1, 10.8]} & 49.2 {\scriptsize[42.3, 58.1]} & 28.0 {\scriptsize[17.6, 38.2]} & 4.8 {\scriptsize[4.6, 5.1]} & 1.5 {\scriptsize[1.1, 1.8]}\\
\textbf{physics + sizes (13)} & 9.1 {\scriptsize[7.3, 11.8]} & 6.7 {\scriptsize[4.5, 9.7]} & 49.7 {\scriptsize[42.9, 58.3]} & 28.0 {\scriptsize[20.7, 36.5]} & 4.1 {\scriptsize[3.7, 4.5]} & 2.6 {\scriptsize[2.1, 3.1]}\\
physics + 141 & 12.7 {\scriptsize[11.4, 14.5]} & 6.4 {\scriptsize[4.0, 9.6]} & 45.5 {\scriptsize[36.5, 55.7]} & 26.3 {\scriptsize[14.8, 37.9]} & 4.8 {\scriptsize[4.5, 5.0]} & 1.5 {\scriptsize[1.1, 1.8]}\\
physics + voltages & \textbf{9.0} {\scriptsize[6.8, 11.9]} & 6.5 {\scriptsize[3.9, 9.9]} & 46.0 {\scriptsize[41.6, 51.5]} & 20.0 {\scriptsize[14.7, 25.7]} & 3.8 {\scriptsize[3.5, 4.2]} & 1.7 {\scriptsize[1.3, 2.0]}\\
signal sizes + voltages & 16.1 {\scriptsize[13.7, 19.5]} & 9.6 {\scriptsize[6.0, 14.6]} & 50.6 {\scriptsize[43.6, 58.3]} & 23.4 {\scriptsize[17.8, 30.0]} & \textbf{3.5} {\scriptsize[3.1, 3.8]} & 1.8 {\scriptsize[1.4, 2.1]}\\
141 + voltages & 16.6 {\scriptsize[15.1, 18.4]} & 6.8 {\scriptsize[3.6, 11.0]} & 48.6 {\scriptsize[44.4, 54.1]} & \textbf{17.7} {\scriptsize[10.4, 25.1]} & 4.4 {\scriptsize[4.1, 4.7]} & \textbf{1.0} {\scriptsize[0.8, 1.2]}\\
\textbf{physics + sizes + voltages} & 9.2 {\scriptsize[7.0, 12.2]} & 6.4 {\scriptsize[3.8, 9.9]} & 44.5 {\scriptsize[40.2, 49.9]} & 18.8 {\scriptsize[13.2, 24.7]} & 3.7 {\scriptsize[3.4, 4.1]} & 1.6 {\scriptsize[1.2, 2.0]}\\
everything & 12.7 {\scriptsize[11.1, 14.9]} & \textbf{6.1} {\scriptsize[3.5, 9.5]} & \textbf{41.1} {\scriptsize[35.5, 47.7]} & 17.8 {\scriptsize[9.9, 25.4]} & 4.3 {\scriptsize[4.0, 4.5]} & 1.1 {\scriptsize[0.8, 1.3]}\\
\bottomrule
\end{tabular}
\end{center}
\paragraph{Reading.}
\begin{itemize}
\item \textbf{Noise, new amplifier}: the physics values carry it (9.1\,\%); signal sizes add nothing to them
  (physics + signal sizes 9.1, + voltages 9.2). The 141 general features are poor alone (16.2) and make the physics values worse
  (12.7): with $\approx155$ training rows they offer cuts that only describe the two training amplifiers.
  Physics + voltages is 9.0 here against 8.6 in test 1, which also has the derived input
  $\ln\sigma_{500}-\widehat{\ln g}(\theta)$ and keeps the three G rows with $\lambda<0.005$; the difference lies well inside the interval.
\item \textbf{Noise, seen amplifier}: every trace set gives 6--7\,\%: the features also identify the amplifier.
\item \textbf{Dark charge}: the only target where more inputs keep helping, 66.8 $\to$ 41.1\,\% (new) and 46.9 $\to$ 17.7\,\% (seen).
  Physics + signal sizes + voltages reaches 44.5 / 18.8 with 44 inputs instead of 185.
\item \textbf{Gain, new amplifier}: no trace input beats the voltages (3.5--3.6\,\%); seen amplifiers reach 1.0\,\% with the 141 features,
  again by identifying the amplifier.
\item Differences below about 2 points for $\sigma_e$ and 5--10 points for $\lambda$ are within the intervals.
\end{itemize}

\paragraph{Summary of the appendix.} Tree models on physics values plus voltages are the most reliable everywhere; the
networks and the deep-kernel GP overfit. Changes are easier than absolute values (the amplifier's own level cancels), which is
what an optimiser needs.

\caveat{Differences between the v3 report's model list and the code: the report says the trace-physics $\lambda$ and gain come from
a ``hierarchical'' fit, the code uses a plain ARD GP (\code{kind="none"}); it says the settings-change models also see the
starting voltages, the code uses $\Delta\theta$ only; the CNNs see all trace curves, not only the noise spectrum; the
pretraining traces are those of the processed directory itself (``pretrained on v2'' could not be checked). The physics tests'
change $\rho$ (all H pairs per amplifier, no drift correction) is not comparable with the scorecard's random-pair $\rho$.}

%=====================================================================================================================
\section{Code map}
\begin{center}\small
\begin{tabularx}{\linewidth}{@{}lX@{}}\toprule
step & code\\\midrule
$\sigma_1(\tau)$, window fit & \code{amp\_phys/trap\_from\_trace.py} (\code{sigmas}, \code{integrals})\\
$\sigma_{500}$, per-packet medians & \code{amp\_phys/physics\_features.py};\newline table \code{amp\_phys/results/per\_channel/physics\_features\_v3.csv}\\
signal sizes & \code{amp\_phys/gain\_features.py}, \code{amp\_phys/clock\_edges.py}\\
8.6\,\% noise forest & \code{amp\_phys/test1\_physics\_model.py}; \code{amp\_phys/results/test1/test1.md}\\
forest on every input type & \code{amp\_phys/ert\_inputs.py}; \code{amp\_phys/results/ert\_inputs/ert\_inputs.csv}\\
labels, usability flags & \code{src/process\_packet.py} (\code{labels\_from\_run\_waders})\\
knobs, classifier, fleet GP, EI, replay & \code{lvl1/fleet\_tests.py}; \code{amp\_phys/results/fleet\_bo/}\\
live search & \code{lvl1/propose\_next.py}; rules \code{configs/limits.toml}, \code{src/ccd\_limits.py}\\
other models & \code{lvl1/scorecard\_v2.py}, \code{lvl1/scorecard.py}, \code{lvl1/train\_lvl1.py}, \code{src/ml/};\newline \code{amp\_phys/test2\_gain\_lambda.py}\\\bottomrule
\end{tabularx}
\end{center}

%=====================================================================================================================
\begin{thebibliography}{99}\small
\bibitem{bonilla2007} E.~V.~Bonilla, K.~M.~A.~Chai, C.~K.~I.~Williams, ``Multi-task Gaussian process prediction'', \emph{Advances in Neural Information Processing Systems} 20 (2008).
\bibitem{breiman1984} L.~Breiman, J.~H.~Friedman, R.~A.~Olshen, C.~J.~Stone, \emph{Classification and Regression Trees} (Wadsworth, 1984).
\bibitem{breiman2001} L.~Breiman, ``Random forests'', \emph{Machine Learning} 45, 5--32 (2001).
\bibitem{byrd1995} R.~H.~Byrd, P.~Lu, J.~Nocedal, C.~Zhu, ``A limited memory algorithm for bound constrained optimization'', \emph{SIAM J.\ Sci.\ Comput.} 16, 1190--1208 (1995).
\bibitem{efron1979} B.~Efron, ``Bootstrap methods: another look at the jackknife'', \emph{Ann.\ Stat.} 7, 1--26 (1979).
\bibitem{elfwing2018} S.~Elfwing, E.~Uchibe, K.~Doya, ``Sigmoid-weighted linear units for neural network function approximation in reinforcement learning'', \emph{Neural Networks} 107, 3--11 (2018).
\bibitem{eriksson2021} D.~Eriksson, M.~Jankowiak, ``High-dimensional Bayesian optimization with sparse axis-aligned subspaces'', \emph{Proc.\ UAI}, PMLR 161, 493--503 (2021).
\bibitem{friedman2001} J.~H.~Friedman, ``Greedy function approximation: a gradient boosting machine'', \emph{Ann.\ Stat.} 29, 1189--1232 (2001).
\bibitem{gardner2014} J.~R.~Gardner, M.~J.~Kusner, Z.~Xu, K.~Q.~Weinberger, J.~P.~Cunningham, ``Bayesian optimization with inequality constraints'', \emph{Proc.\ ICML}, PMLR 32, 937--945 (2014).
\bibitem{gelbart2014} M.~A.~Gelbart, J.~Snoek, R.~P.~Adams, ``Bayesian optimization with unknown constraints'', \emph{Proc.\ UAI} (2014).
\bibitem{geurts2006} P.~Geurts, D.~Ernst, L.~Wehenkel, ``Extremely randomized trees'', \emph{Machine Learning} 63, 3--42 (2006).
\bibitem{hanley1982} J.~A.~Hanley, B.~J.~McNeil, ``The meaning and use of the area under a receiver operating characteristic (ROC) curve'', \emph{Radiology} 143, 29--36 (1982).
\bibitem{hanley2006} J.~A.~Hanley, B.~J.~MacGibbon, ``Creating non-parametric bootstrap samples using Poisson frequencies'', \emph{Comput.\ Methods Programs Biomed.} 83, 57--62 (2006).
\bibitem{hoffman2014} M.~D.~Hoffman, A.~Gelman, ``The No-U-Turn sampler'', \emph{J.\ Mach.\ Learn.\ Res.} 15, 1593--1623 (2014).
\bibitem{hopkinson1982} G.~R.~Hopkinson, D.~H.~Lumb, ``Noise reduction techniques for CCD image sensors'', \emph{J.\ Phys.\ E} 15, 1214--1222 (1982).
\bibitem{huber1964} P.~J.~Huber, ``Robust estimation of a location parameter'', \emph{Ann.\ Math.\ Stat.} 35, 73--101 (1964).
\bibitem{jones1998} D.~R.~Jones, M.~Schonlau, W.~J.~Welch, ``Efficient global optimization of expensive black-box functions'', \emph{J.\ Global Optim.} 13, 455--492 (1998).
\bibitem{ke2017} G.~Ke \emph{et al.}, ``LightGBM: a highly efficient gradient boosting decision tree'', \emph{Advances in Neural Information Processing Systems} 30 (2017).
\bibitem{koenker1978} R.~Koenker, G.~Bassett, ``Regression quantiles'', \emph{Econometrica} 46, 33--50 (1978).
\bibitem{lakshminarayanan2017} B.~Lakshminarayanan, A.~Pritzel, C.~Blundell, ``Simple and scalable predictive uncertainty estimation using deep ensembles'', \emph{Advances in Neural Information Processing Systems} 30 (2017).
\bibitem{loshchilov2019} I.~Loshchilov, F.~Hutter, ``Decoupled weight decay regularization'', \emph{Proc.\ ICLR} (2019).
\bibitem{nix1994} D.~A.~Nix, A.~S.~Weigend, ``Estimating the mean and variance of the target probability distribution'', \emph{Proc.\ IEEE Int.\ Conf.\ Neural Networks} (1994).
\bibitem{oord2018} A.~van den Oord, Y.~Li, O.~Vinyals, ``Representation learning with contrastive predictive coding'', arXiv:1807.03748 (2018).
\bibitem{owen1998} A.~B.~Owen, ``Scrambling Sobol' and Niederreiter--Xing points'', \emph{J.\ Complexity} 14, 466--489 (1998).
\bibitem{pearson1901} K.~Pearson, ``On lines and planes of closest fit to systems of points in space'', \emph{Phil.\ Mag.} 2, 559--572 (1901).
\bibitem{pedregosa2011} F.~Pedregosa \emph{et al.}, ``Scikit-learn: machine learning in Python'', \emph{J.\ Mach.\ Learn.\ Res.} 12, 2825--2830 (2011).
\bibitem{rasmussen2006} C.~E.~Rasmussen, C.~K.~I.~Williams, \emph{Gaussian Processes for Machine Learning} (MIT Press, 2006).
\bibitem{schonlau1998} M.~Schonlau, W.~J.~Welch, D.~R.~Jones, ``Global versus local search in constrained optimization of computer models'', in \emph{New Developments and Applications in Experimental Design}, IMS Lecture Notes 34, 11--25 (1998).
\bibitem{sobol1967} I.~M.~Sobol', ``On the distribution of points in a cube and the approximate evaluation of integrals'', \emph{USSR Comput.\ Math.\ Math.\ Phys.} 7(4), 86--112 (1967).
\bibitem{spearman1904} C.~Spearman, ``The proof and measurement of association between two things'', \emph{Am.\ J.\ Psychol.} 15, 72--101 (1904).
\bibitem{stein1999} M.~L.~Stein, \emph{Interpolation of Spatial Data: Some Theory for Kriging} (Springer, 1999).
\bibitem{swersky2013} K.~Swersky, J.~Snoek, R.~P.~Adams, ``Multi-task Bayesian optimization'', \emph{Advances in Neural Information Processing Systems} 26 (2013).
\bibitem{tiffenberg2017} J.~Tiffenberg \emph{et al.}, ``Single-electron and single-photon sensitivity with a silicon Skipper CCD'', \emph{Phys.\ Rev.\ Lett.} 119, 131802 (2017).
\bibitem{wilson2016} A.~G.~Wilson, Z.~Hu, R.~Salakhutdinov, E.~P.~Xing, ``Deep kernel learning'', \emph{Proc.\ AISTATS}, PMLR 51, 370--378 (2016).
\end{thebibliography}

\end{document}
